\section{Introduction}
Simultaneous speech translation (SimulST) systems must produce target text incrementally as source speech streams in, making latency as critical as translation quality. Unlike offline translation, where the full utterance is available before decoding begins, SimulST must decide at each step whether to wait for more input or commit to a partial translation: a decision with lasting consequences, since committed tokens cannot be revised.

Latency in SimulST arises from many sources --- speech segmentation, ASR delay, model decoding, and policy instability among them --- but two are particularly tied to the read/write decision itself. The first is structural: language pairs differ in canonical constituent order and head direction, so English-to-Japanese requires verb-final reordering, while English-to-Chinese exhibits phrase-level head-direction and discourse-order divergences that still force the translator to delay output until enough source words arrive. The second is semantic: even when word order is not an issue, the correct translation of a current segment may depend on future context not yet heard --- a pronoun whose referent appears later, a verb whose argument structure changes the meaning of an earlier noun. Existing methods tend to address these in isolation. Word alignment-based approaches such as InfiniSST explicitly model structural reordering delays but have no mechanism to reason about semantic future dependency. LLM-based synthesis methods such as EAST leverage broader contextual knowledge but provide no quality guarantees and exhibit unstable latency behavior under both conditions. More fundamentally, existing trajectory-synthesis pipelines commit on the basis of the observed prefix or of literal prefix agreement across hypotheses, treating commitment as a question about what has already been seen.

We argue that safe commitment is fundamentally a \emph{counterfactual} question: would this target prefix remain valid under plausible futures of the source? We propose [METHOD], a training-data synthesis framework that operationalizes this view as \emph{distributional consensus over future-conditioned next-token distributions}. Given a partial source utterance, our method samples multiple plausible future source continuations from a base language model. For each hypothesized full source, a larger instruction-tuned LLM produces the next-token distribution over the target vocabulary conditioned on the committed translation prefix. The method then selects the token that lies in the intersection of high-probability candidate sets across every future, committing only tokens on which all futures agree. This token-level consensus loop repeats until no intersection exists, naturally yielding the longest committable prefix that is robust to future uncertainty. The resulting consensus translations, paired with their source prefixes and commitment points, serve as training data for a SimulST model --- distilling this expensive future-conditioned reasoning into a policy whose inference-time cost is unchanged, without requiring reinforcement learning or inference-time LLM augmentation.

We evaluate on GigaSpeech and RealSI across English$\rightarrow$Chinese, English$\rightarrow$German, and English$\rightarrow$Japanese, comparing against strong baselines including EAST, Simul-MuST-C, LA-$N$, wait-$k$, InfiniSST, and TAF along the quality--latency tradeoff curve. Our goal is to test whether grounding commitment in distributional future-conditioned consensus offers a more principled and cost-effective basis for SimulST data synthesis than prefix-only alternatives.